
\textbf{Corollaire 4.15.} \uline{Soit} \(B'\) \uline{le} \(\lambda\)-\uline{anneau engendré sur} \(K\) \uline{par l'indéterminée} \(X\) \uline{soumise aux relations} \(\gamma^k(X)=0\) \uline{pour} \(k>n\), \uline{i}.\uline{e}. \(B'=K[X,\ldots,\gamma^n X]\) (4.9). \uline{Alors la filtration de} \(\lambda\)-\uline{anneau de} \(B'\) \uline{est identique à la filtration par le poids par rapport aux} \(\gamma^j X\).

Compte tenu de la proposition 4.12, de la remarque 4.13.ii) et du fait que \(A=K\), \(\Fil^k B'=K\mathbin{.} M_k+B'\mathbin{.} M_{k+1}\), donc \(\Fil^k B'\) est l'idéal engendré par les monômes de poids \(\ge k\) par rapport aux \(\gamma^j X\).

\section*{5. Les \(\lambda^p(N,x)\).}

Ce paragraphe ne sera utilisé que dans Exp.VII~4, pour l'étude des immersions régulières~; il peut donc être omis en première lecture.

5.1. Soit \(d\) un entier \(\ge1\). On désigne par \(A\) la \(\mathbb Z\)-\(\lambda\)-algèbre engendrée par les générateurs \(N\), \(x\) soumis aux relations \(\lambda^i(N)=0\) pour \(i>d\)~; on a donc :
\[
  A=\mathbb Z[N,\ldots,\lambda^dN;
       x,\ldots,\lambda^jx,\ldots].
\]

On pose \(\lambda_{-1}(N)=\sum_{i=0}^{d}(-1)^i\lambda^i(N)\).

\textbf{Lemme 5.2.} \uline{Pour tout} \(p\ge1\), \(\lambda^p(x\mathbin{.}\lambda_{-1}(N))\) \uline{est divisible par} \(\lambda_{-1}(N)\).

On sait (2.3) que \(\lambda^p(x\mathbin{.}\lambda_{-1}(N))=P_p(\lambda^i x,\lambda^j(\lambda_{-1}(N)))\) est un polynôme isobare de poids \(p\) par rapport aux \(\lambda^j(\lambda_{-1}(N))\)~; il suffit donc de montrer le lemme pour les \(\lambda^p(\lambda_{-1}(N))\).

Soit \(B\) le \(\lambda\)-anneau engendré sur \(\mathbb Z\) par les indéterminées \(N_1,\ldots,N_d\), soumises aux relations \(\lambda^i(N_j)=0\) si \(i\ge2\). Il est clair que \(\lambda^k(\sum_{i=1}^{d}N_i)=0\) pour \(k>d\). On peut donc définir un \(\lambda\)-homomorphisme \(f\) de \(A\) dans \(B\), en envoyant \(N\) sur \(\sum_{i=1}^{d}N_i\)~; on a alors \(\Ghat(f)(\lambda_t(N))=\prod_{i=1}^{d}\lambda_t(N_i)=\prod_{i=1}^{d}(1+N_it)\),
